\section{Overview}
\label{sec:intro_overview}

Autonomous perception is one of the most critical frontiers in automotive engineering, driven by consumer demand and regulatory mandates for \acrfull{adas}.
However, the pursuit of full autonomy has driven a relentless escalation in the complexity and resolution of these perception systems:
modern vehicles rely on \acrfullpl{cnn} to process high-definition video streams in real-time,
executing billions of dense, \acrfullpl{flop} every second.

As the industry scales towards higher resolutions and faster reaction times to ensure passenger safety, the energy demands of this continuous, dense spatial reasoning have grown substantially.
\textcite{lin_adas_cnn_issues} documented CPU-based \gls{cnn} inference exceeding the thermal and power constraints of \gls{adas} hardware, motivating the industry's shift to \gls{gpu}-class accelerators; 
continued growth in perception complexity now places dense \gls{ann} workloads on \glspl{gpu} on a similar trajectory, motivating a paradigm shift towards fundamentally more efficient architectures.

This project investigates utilizing cutting-edge neuromorphic paradigms, more specifically \acrfullpl{snn},
and their feasibility in this domain as an alternative computational approach.
To achieve this, the project develops \texttt{DetectorNX-G3-SNN}, a high-definition spiking architecture designed for multi-object automotive perception. By rigorously benchmarking this system against a strictly equivalent continuous \gls{cnn} on the event-based \acrfull{etram} dataset \parencite{Verma_2024_CVPR}, this research empirically deconstructs the performance-efficiency frontier---measuring the precise trade-offs between the energy savings of event-driven dynamics and the spatial precision constraints of 1-bit temporal quantization.